\section{Conclusion}\label{sec:conclusion}

In a linear market with feedback traders, prices summarize the trader population in one object, the feedback kernel. The kernel is sufficient for every outcome that price-based trading determines, for any insider and any causal dealer rule. It is all that prices identify, and with inventory-averse dealers both statements hold in equilibrium with unchanged memories. Information beyond the kernel, who the groups are and how much each trades, is worth nothing to anyone trading on public information. It is worth something to a supervisor whose action is aimed at a particular group. Learning it from prices costs of order $(\Delta\eta)^{-(4K-2)}$ experiments or observations, executed quantity of order $\delta^{-2}$ and capital of order $\delta^{-4}$. Class-tagged flow removes that cost. With position limits, a second layer appears: how close traders are to their limits, which calm-market prices do not reveal at all under hard limits, and which disclosure of capacities or stress alone can show.

The practical reading is a division of supervisory labour. Market-wide measures, such as limits on aggregate feedback intensity or confidence sets for the price response, can be built from public data, because they depend on the kernel alone. Targeted measures cannot be built from public data at reasonable cost when the groups' memories are close. For targeted measures, verified class-level records are the efficient source of information, not more price data or more sophisticated estimation.
